\chapter{Conclusions and Future Work}

\section{Conclusions}

This thesis started as a dense deep-learning surrogate for UAV radio maps and
ended as a hybrid prior-anchored CKM prediction pipeline. The main lesson is
that architecture alone was not enough. The decisive improvements came from
understanding the target distributions, enforcing strict ground-only metrics,
separating LoS and NLoS behaviour, and using physical priors where the dataset
clearly followed physical structure.

The original targets were met by the finished non-DL priors for path loss,
delay spread, and angular spread. Try 78 achieved \SI{1.9246}{dB} overall
path-loss RMSE, with \SI{1.7516}{dB} LoS and \SI{3.3967}{dB} NLoS. Try 79
achieved \SI{28.10}{ns} delay-spread RMSE and \SI{13.74}{\degree}
angular-spread RMSE on the test split. The final Try 80 neural model is still
training; its final test results must be inserted once the checkpoint is
selected.

The main contribution is the final synthesis: a CKM pipeline that combines
calibrated propagation priors, city-holdout discipline, sinusoidal UAV-height
conditioning, topology/regime diagnostics, and distribution-first residual
heads. This is different from both a pure empirical propagation formula and a
pure black-box CNN.

\section{Future Directions}

The most important future extension is frequency conditioning. The current
dataset fixes the radio configuration, while a general outdoor CKM model should
learn how path loss and spreads vary with carrier frequency, bandwidth, antenna
pattern, and possibly polarization. A natural next dataset would therefore
contain the same cities simulated at multiple frequencies and antenna settings.

The second extension is field calibration. The priors in Try 78 and Try 79 are
already interpretable, so a small number of real measurements could be used to
calibrate the simulator-to-reality gap without retraining the full model from
scratch.

Finally, the Try 80 idea can be expanded with stronger global-context
backbones, transmitter-oriented mixup, full corridor losses, and test-time
adaptation. These should be added only after preserving the current safeguards:
strict city holdout, ground-only masking, and explicit LoS/NLoS reporting.
